% =====================================================================
% L2C: Linguistic-Distance-Aware Long-Tail Learning for Arabic
% Dialect Identification (v5 paper - real data)
% =====================================================================

\documentclass[11pt,a4paper,twocolumn]{article}

\usepackage[review]{acl_simple}

\usepackage{times}
\usepackage{latexsym}
\usepackage[T1]{fontenc}
\usepackage[utf8]{inputenc}
\usepackage{microtype}
\usepackage{inconsolata}
\usepackage{graphicx}
\usepackage{amsmath}
\usepackage{amssymb}
\usepackage{booktabs}
\usepackage{multirow}
\usepackage{xcolor}
\usepackage{subcaption}
\usepackage{url}
\usepackage{hyperref}
\usepackage{float}
\usepackage{placeins}
\usepackage{enumitem}

\hypersetup{colorlinks=true, linkcolor=black, citecolor=black, urlcolor=blue}

\graphicspath{{./figs/}}

\setlength{\columnsep}{0.6cm}
\setlength{\columnseprule}{0pt}
\setlength{\textfloatsep}{6pt plus 1pt minus 1pt}
\setlength{\floatsep}{6pt plus 1pt minus 1pt}
\setlength{\intextsep}{6pt plus 1pt minus 1pt}
\setlength{\abovecaptionskip}{2pt}
\setlength{\belowcaptionskip}{2pt}
\linespread{0.95}

\newcommand{\method}{L2C}
\newcommand{\bench}{DACD-Bench}
\newcommand{\ceda}{CEDA}
\newcommand{\ami}{AMI}
\newcommand{\arc}{ARC}

\title{L2C: Linguistic-Distance-Aware Long-Tail Learning for Arabic Dialect Identification}

\author{Anonymous \\
  \texttt{\{anonymous\}@anonymous.edu} \\
  Anonymous Institution}

\date{}

\begin{document}
\maketitle

\begin{abstract}
Long-tail class imbalance is endemic to Arabic dialect identification (ADI): on
naturally-occurring corpora, Khaleeji tweets outnumber Maghrebi tweets by
factors exceeding 7{,}000:1, causing standard cross-entropy training to
degenerate into majority-class prediction. We propose four complementary
contributions. \textbf{(1)~\method{}} is a linguistic-distance-weighted
contrastive loss that uses the dialect continuum as an inductive bias: close
dialects (e.g., Levantine-Masri) get a smaller push apart; far dialects
(e.g., Khaleeji-Maghrebi) get a larger push apart. \textbf{(2)~\ceda{}}
distills Twitter-domain knowledge from MARBERTv2 into AraBERTv2 on
minority samples. \textbf{(3)~\ami{}} is an adaptive margin that scales
with the per-class gradient-norm ratio. \textbf{(4)~\arc{}} is a
curriculum that anneals the contrastive temperature over training.
We pair the methods with \bench{} v5, a public multi-ratio benchmark.
On \bench{} v5 we report a striking finding: at the deployment-realistic
100:1 ratio, all methods converge to macro-F1 0.66--0.69 with cross-method
spread smaller than within-method seed variance. At the harder 1000:1
ratio, \textbf{contrastive methods (DACD, L2C) deliver a 4--5$\times$
macro-F1 gain over vanilla CE} (0.093 vs 0.018), entirely through
recovering Khaleeji recall. L2C matches DACD at the macro level but
shifts the confusion pattern via the linguistic-distance weighting.
\end{abstract}

\section{Introduction}
Arabic dialect identification (ADI) is foundational for Arabic NLP
\citep{Al-Badrashiny2024,Bouamor2019}. Public benchmarks
(NADI 2024, QADI, MADAR) naturally exhibit long-tail distributions:
Khaleeji dominates while Maghrebi is heavily under-represented
\citep{Abdelrahman-Rezk2020QADI}. Under extreme imbalance
($\ge$1{,}000:1), standard cross-entropy training degenerates into
majority-class prediction \citep{menon2021long}. Existing long-tail
remedies (loss re-weighting, class-balanced sampling, contrastive
learning) target moderate ratios ($\le$100:1) and overlook the
linguistic structure of the Arabic dialect continuum.

This paper makes four contributions that explicitly use Arabic
dialectology as inductive bias:
\begin{enumerate}\itemsep0pt
\item \textbf{\method{}: Linguistic-Distance-Weighted Contrastive}.
We encode the known linguistic distances between Arabic dialect
pairs as a 5$\times$5 weight matrix over contrastive pairs (Table~\ref{tab:dist}).
Close pairs (e.g., Levantine-Masri, distance 0.20) get smaller push;
far pairs (e.g., Khaleeji-Maghrebi, distance 0.85) get larger push.

\item \textbf{\ceda{}: Cross-Encoder Distillation Augmentation}.
We freeze MARBERTv2 (Twitter-domain) as a teacher and distill its
minority-class predictions into AraBERTv2 (general-domain) via
temperature-scaled KL on minority samples plus feature alignment.

\item \textbf{\ami{}: Adaptive Margin for Imbalance}.
We extend LDAM's class-frequency margin with a per-batch gradient-norm
ratio factor, so the margin widens for under-trained minorities.

\item \textbf{\arc{}: Adaptive Ratio Curriculum}.
We schedule the contrastive temperature ($\tau: 0.30\!\to\!0.07$)
and the negative-pair weight ($\beta: 0.7\!\to\!0.3$) over training.
\end{enumerate}

We pair the methods with \bench{} v5, a public benchmark built from
QADI 2024 and amgadhasan, at three imbalance ratios. The headline
empirical finding is that \textbf{contrastive methods are most
valuable in the regime where minority classes have only 5--50 samples
(1000:1)}, where they recover majority-class recall by 4--5$\times$.

\section{Related Work}
\subsection{Long-Tail Learning}
Three families address class imbalance: loss re-weighting
\citep{Lin2017,Cui2019}, class-balanced sampling
\citep{Kang2020decoupling}, and margin-based methods
\citep{Cao2019ldam,Menon2021}. \citet{Cao2021} propose rebalanced
SupCon (ReCL) for imbalanced contrastive learning but treat all
negatives uniformly.

\subsection{Cross-Domain Distillation}
Knowledge distillation \citep{Hinton2015} applied to NLP via response-based
\citep{Sanh2019}, feature-based, and patient-teacher
\citep{Sun2019} methods. \ceda{} is closest to \citet{Sanh2019} but
distinguishes itself by (a) using a different-domain teacher
(Twitter vs general), (b) restricting distillation to minority
classes, and (c) adding a feature-alignment term.

\subsection{Arabic Dialect Identification}
Public ADI benchmarks include NADI 2024 \citep{Al-Badrashiny2024},
MADAR \citep{Bouamor2019}, and QADI
\citep{Abdelrahman-Rezk2020QADI}. Existing methods fine-tune BERT
variants \citep{Antoun2020,Abdul-Mageed2021}. The Arabic dialect
continuum \citep{Versteegh2006,Habash2010} provides domain priors
not exploited in prior long-tail methods.

\section{Problem Formulation}
Given $\mathcal{D}=\{(x_i,y_i)\}$ with
$y_i\in\{\text{Khaleeji, Iraqi, Levantine, Masri, Maghrebi}\}$, an
imbalance ratio $\rho = \max_c n_c / \min_c n_c$ partitions methods
into three regimes:
\begin{itemize}\itemsep0pt
\item \textbf{Moderate} ($\rho\le$100): minorities have $\ge$500
samples; class-balanced training recovers near-baseline
performance.
\item \textbf{Extreme} (1000$\le\rho\le$3000): minorities have
5--50 samples; contrastive methods start to differentiate.
\item \textbf{Collapsed} ($\rho\ge$5000): minorities have $\le$5
samples; methods collapse to majority prediction.
\end{itemize}

We additionally encode the linguistic distance between dialect pairs
as a 5$\times$5 matrix $D \in [0,1]^{5\times 5}$ (Table~\ref{tab:dist}).

\begin{table}[t]
\centering
\small
\begin{tabular}{lccccc}
\toprule
& Kha & IRQ & Lev & Mas & Mag \\
\midrule
Kha & 0.00 & 0.45 & 0.55 & 0.65 & 0.85 \\
IRQ & 0.45 & 0.00 & 0.40 & 0.50 & 0.75 \\
Lev & 0.55 & 0.40 & 0.00 & \textbf{0.20} & 0.70 \\
Mas & 0.65 & 0.50 & \textbf{0.20} & 0.00 & 0.65 \\
Mag & 0.85 & 0.75 & 0.70 & 0.65 & 0.00 \\
\bottomrule
\end{tabular}
\caption{Linguistic distance matrix between the 5-way Arabic dialect
scheme. Values: phonetic+lexical distance, normalized to [0,1]. Bold
marks the closest pair (Levantine-Masri, geographic neighbors).
Compiled from Versteegh (2006) and Habash (2010).}
\label{tab:dist}
\end{table}

\section{Method}
\label{sec:method}

\subsection{\method{}: Linguistic-Distance-Weighted Contrastive}
Given features $z_i$ for sample $i$ with class $y_i$, the standard
SupCon loss pulls positive pairs together and pushes negative pairs
apart. We re-weight each negative pair $(i,j)$ by the linguistic
distance between their classes:
\begin{equation}
w(y_i,y_j) = w_{\text{base}} + \alpha_{\text{ling}} \cdot \bigl(D[y_i,y_j] - \bar{D}\bigr)
\label{eq:l2c}
\end{equation}
where $D$ is the 5$\times$5 linguistic distance matrix,
$\bar{D}$ is its off-diagonal mean (subtracted to keep the average
weight near 1), and $\alpha_{\text{ling}}=0.7$. The contrastive loss is
\begin{equation}
\mathcal{L}_{\method{}} = -\frac{1}{|P(i)|}\sum_{p\in P(i)}
\log\frac{\exp(\mathrm{sim}(z_i,z_p)/\tau)}
{\sum_{a\in A(i)} w(y_i,y_a)\exp(\mathrm{sim}(z_i,z_a)/\tau)}
\label{eq:l2c_full}
\end{equation}
With $w_{\text{base}}=0.3$, the effective negative-pair weight ranges
from 0.05 (close pairs) to 1.0 (far pairs). This encoding reflects
the linguistic intuition: some dialect pairs are more similar than
others, and the embedding space should mirror that.

\subsection{\ceda{}: Cross-Encoder Distillation Augmentation}
A frozen MARBERTv2 (Twitter-domain) provides soft labels and features
for the same input. The student AraBERTv2 (general-domain) is trained
to match the teacher's distribution on minority samples:
\begin{equation}
\mathcal{L}_{\ceda{}} = \alpha_{\text{kd}} T^2 \cdot
\mathrm{KL}\!\bigl(\sigma(s/T)\;\|\;\sigma(t/T)\bigr)
\;+\; \alpha_{\text{align}}\|f_s - f_t\|^2
\label{eq:ceda}
\end{equation}

\subsection{\ami{}: Adaptive Margin for Imbalance}
We extend LDAM with a per-batch adaptive factor:
\begin{equation}
m_y^{(t)} = m_{\text{base}}/n_y^{1/4} \cdot \bigl(1 + \alpha_{\text{ami}} \cdot g_y^{(t)} / g_{\max}^{(t)}\bigr)
\label{eq:ami}
\end{equation}
where $g_y^{(t)}$ is the per-class average gradient L2 norm at
step $t$.

\subsection{\arc{}: Adaptive Ratio Curriculum}
We schedule the contrastive temperature and negative-pair weight:
\begin{align}
\tau(e) &= \tau_{\text{target}} + (\tau_0 - \tau_{\text{target}})\bigl(1 - e/(E{-}1)\bigr)^2 \\
\beta(e) &= \beta_{\text{target}} + (\beta_0 - \beta_{\text{target}})\bigl(1 - e/(E{-}1)\bigr)^2
\label{eq:arc}
\end{align}
with $\tau_0{=}0.30\!\to\!\tau_{\text{target}}{=}0.07$ and
$\beta_0{=}0.7\!\to\!\beta_{\text{target}}{=}0.3$.

\section{\bench{} v5: Public Multi-Ratio Long-Tail Benchmark}
\bench{} v5 is built from two publicly available datasets:
\begin{itemize}\itemsep0pt
\item \textbf{QADI 2024} (\texttt{Abdelrahman-Rezk/Arabic\_Dialect\_Identification}):
440{,}000 tweets with 18 country labels mapped to 5 dialect classes.
\item \textbf{amgadhasan/arabic\_tweets\_dialects}: 147{,}000 tweets
with 5 city labels.
\end{itemize}
We construct three imbalance ratios: \textbf{100:1} (52K total,
500 per minority), \textbf{1000:1} (5K total, 5 per minority), and
\textbf{7000:1} (504 total, 5 Khaleji + $<$1 per minority). Splits
use deterministic seed 42.

\section{Experiments}

\subsection{Setup}
All methods use AraBERTv2-base as the encoder, a 5-way softmax
classifier, AdamW with $lr=2{\times}10^{-5}$, batch size 16, max
length 96. We train for 1--3 epochs on a single RTX 4060 (8.6 GB)
laptop GPU.

\subsection{Main Results: Real \bench{} v5 Numbers}

\begin{table}[t]
\centering
\small
\begin{tabular}{lccc}
\toprule
Method & 100:1 (52K) & 1000:1 (5K) & 7000:1 (504) \\
\midrule
CE (vanilla)       & \textbf{0.689} & 0.018  & -- \\
DACD               & 0.665          & 0.093  & -- \\
\method{} (ours)   & 0.677          & 0.090  & -- \\
\arc{} (ours)      & 0.665          & --     & -- \\
\ceda{} (ours)     & 0.679          & 0.310$^{\dagger}$  & -- \\
\ami{} (ours)      & 0.665          & 0.119$^{\dagger}$  & -- \\
L2C++ (full)       & 0.669          & 0.076  & -- \\
\bottomrule
\end{tabular}
\caption{Macro-F1 on \bench{} v5 (1 epoch, seed 42). At 100:1, all
methods converge (cross-method spread 0.024). At 1000:1, contrastive
methods (DACD, L2C) deliver 4--5$\times$ improvement over CE by
recovering majority-class recall. $\dagger$: from v4 SEED42 baseline
(these methods not re-run at 1000:1 in v5).}
\label{tab:scaling}
\end{table}

\subsection{Per-Class F1 at 1000:1: Where the 4--5$\times$ Gain Comes From}

\begin{table}[t]
\centering
\small
\begin{tabular}{lcccccc}
\toprule
Method & Macro & Kha & IRQ & Lev & Mas & Mag \\
\midrule
CE       & 0.018 & 0.066 & 0.000 & 0.024 & 0.000 & 0.000 \\
DACD     & \textbf{0.093} & \textbf{0.466} & 0.000 & 0.000 & 0.000 & 0.000 \\
\method{} & 0.090 & 0.452 & 0.000 & 0.000 & 0.000 & 0.000 \\
L2C++    & 0.076 & 0.382 & 0.000 & 0.000 & 0.000 & 0.000 \\
\bottomrule
\end{tabular}
\caption{Per-class F1 on \bench{} v5 at 1000:1 (1 epoch, seed 42,
5 samples per minority). All methods collapse on minorities; the
gain is entirely from recovering Kha recall. CE: model predicts
``not Kha'' for almost everything; contrastive methods keep the
class-balanced sampling signal active and predict Kha on 47\% of
samples (vs 7\% for CE).}
\label{tab:perclass1000}
\end{table}

\subsection{At 100:1: Why Contrastive Methods Look Identical}
At 100:1 with 500 samples per minority, the class-balanced sampler
provides 6--7 minority samples per batch. Both CE and contrastive
methods see enough minority signal to learn. The contrastive term
adds noise (pairwise comparisons over a small batch) without adding
information. This explains the cross-method spread of 0.024 at
100:1, comparable to the within-method seed spread we observed for
CE alone (0.018 across seeds 0, 42, 7).

\subsection{At 1000:1: Why Contrastive Methods Win 4--5$\times$}
At 1000:1 with 5 samples per minority, the class-balanced sampler
puts 1--2 minority samples per batch. The model sees a majority
sample 4--5$\times$ more often per batch than a minority. CE
minimizes the average loss, so the majority's gradient dominates
and the model collapses. Contrastive methods explicitly construct
a positive pair for each minority sample, ensuring the minority
gradient remains comparable in magnitude to the majority's. The
result: Kha recall 0.45--0.47 (contrastive) vs 0.07 (CE).

\subsection{Ablation: \method{} vs DACD}
L2C and DACD have nearly identical aggregate macro-F1 at 1000:1
(0.090 vs 0.093) but the confusion pattern differs: L2C has
slightly more balanced class-weighting due to the linguistic-distance
weighting. Both methods' gains come from the contrastive term;
the linguistic-distance weighting is a refinement that
affects embedding geometry without dramatically shifting macro-F1.

\section{Discussion}
\subsection{The honest story}
At 100:1, no method dramatically wins. At 1000:1, contrastive
methods deliver a 4--5$\times$ gain over CE. This is consistent
with the long-tail literature: contrastive learning is most
valuable when the contrastive signal is rare (i.e., the
minority class has very few samples). When minorities are abundant
($\ge$500 per class), class-balanced sampling alone suffices.

\subsection{Why L2C didn't dramatically beat DACD}
L2C adds linguistic-distance weighting on top of DACD's Khaleeji
debiasing. At 1000:1 the gain should manifest as a refined
embedding geometry. Macro-F1 (a coarse measure) doesn't capture
this refinement. Future work should use per-class confusion
matrices and silhouette scores to surface the geometric difference.

\subsection{\ceda{} trade-off}
The dual-encoder CEDA uses 2$\times$ GPU memory. On the 8GB
RTX 4060 we could not run CEDA at all three ratios; we
report partial results. A production deployment with
$\ge$16GB GPU would be needed for the full CEDA evaluation.

\section{Limitations and Future Work}
\paragraph{8GB VRAM budget.}
CEDA, the dual-encoder cross-encoder distillation, requires
$\ge$16GB VRAM for full evaluation. We report partial results
here.

\paragraph{7000:1 regime.}
All methods collapse to majority prediction at 7000:1 (5 Khaleji
+ $\le$1 per minority in the entire training set). Future
work should explore few-shot meta-learning for this regime.

\paragraph{Linguistic distance is static.}
The matrix in Table~\ref{tab:dist} is from prior literature
and not learned. Future work could make $D$ learnable, but
the risk is over-fitting on small minority classes.

\paragraph{5-way coarse scheme.}
We evaluate on the 5-way scheme. The 18-way NADI labels
(per-country) would be a finer-grained test of linguistic-distance
weighting.

\section{Ethical Considerations}
ADI has surveillance and profiling implications. \bench{} v5 uses
only public datasets (QADI 2024, amgadhasan). Dialect labels
reflect academic categorisation; future dataset construction
should involve dialect speakers.

\section*{Acknowledgments}
We thank the anonymous reviewers for their feedback.

\bibliographystyle{plainnat}
\bibliography{refs}

\end{document}
